\section{Conclusion}
\label{sec:conclusion}

\subsection{Findings}

\begin{enumerate}[itemsep=2pt]
  \item The DTaaS client has no spatial visualisation capability today beyond embedded
        time series.
  \item Every surveyed commercial platform shares one architecture --- anchor, encode,
        sustain --- and sells a no-code editor over the encode layer.
  \item The renderer is therefore the wrong thing to standardise on. Transport
        normalisation, the temporal store, anchor resolution and the encoding vocabulary
        are the durable artefacts, and none exist off the shelf.
  \item Not every substrate is a 3D scene. Photographs and live video are first-class
        places to put simulation output; generalising the adapter interface costs
        nothing and changes what ships in the first two months.
  \item Roughly a third of the layer is already deployed --- historian, broker, device
        registry, charting, versioned storage and identity need to be \emph{used}, not
        built.
  \item Hospitals and buildings are one 3D problem, machines another, city/mobility/marine
        a third: three geometric adapters plus \texttt{image}, \texttt{video} and
        \texttt{embed}.
  \item Simulation--measurement co-visualisation is open ground, and DTaaS is unusually
        placed to take it.
  \item One early decision governs half the catalogue: every encoding must read
        \texttt{valueAt(playhead)}, never ``latest value''.
  \item Licence discipline removes two of the strongest candidates
        (\pkg{xeokit-sdk}, \pkg{realvirtual WEB}) and flags two more.
\end{enumerate}

\subsection{Positioning and next step}

The literature converges on this architecture from several directions.
\citet{rundel2023leveraging} show that the escape from domain-rebuild cost is to derive
the scene from data that already exists --- and a photograph is data that already exists.
\citet{vasilijevic2024trondheim} adopt a deliberately two-tier visualisation, Grafana for
series and a game-engine/Cesium layer for geospatial rendering, which is almost exactly
the split proposed here between the \texttt{embed} adapter and the geometric substrates,
coupled through a shared clock rather than merged. \citet{ali2024spatial} caution that
\tech{10} is not merely engineering. \citet{mobilitydt2023} supply the evaluation
constructs and, by measuring how expensive geometry reconstruction is, the strongest
argument for using imagery where it suffices. \citet{liao2025digital} demonstrate the
ghost twin in miniature; \citet{chen2024tangible} supply the human-factors constructs and
the case for shared, multi-operator views; and \citet{yang2024digital} identify data
integration and quality, not rendering, as the binding constraint on DT adoption ---
consistent with the claim that the binding layer is the product.

\medskip
\noindent\textbf{Next step.} Build Phase~0 against the services already running, then
Phase~1 on a photograph. Six weeks to something a user recognises as their own twin, no
WebGL required, and the whole binding layer proven before a single 3D asset is
commissioned.
